\documentclass[11pt,a4paper]{ctexart}
\usepackage[color]{tlnote}
\title{高数 I（A 类）期中 · 11/10 执行卡}
\author{陶理}
\noteheader{考前最后做题日（D-1）}
\docdate{2026-11-10}

\begin{document}
\renewcommand{\arraystretch}{1.18}
\noindent{\Large\bfseries 高数 I（A 类）期中 · 11/10 执行卡（D-1）}\hfill{\small\color{tlblue}2026-11-10 / 周三 11-11 考试}
\par\noindent\rule{\linewidth}{0.4pt}\par\smallskip
\noindent{\small\color{gray}目标：明天（11/11，周三）考试。今天是最后一个做题日——收尾两套整卷手感，按诊断卷实测的薄弱点做定向加练 72 题；23:00 停笔，进入纯保温。\textbf{诊断卷 D1–D12 已于昨晚（11/9 19:30）做完}，今天直接进入 72 题，无需重做诊断。}

\begin{keypoint}[今日三块]
\begin{tabular}[t]{@{}p{3.0cm}p{12.0cm}@{}}
早段 & G12 模拟卷 C 18 题（单题回炉，不掐表）＋ 模拟卷 B 全卷 120 min 掐表 ＋ 50 min 订正 \\[2pt]
下午/晚间 & 定向加练 72 题（六模块，A 级优先：M1→M2→M3→M4，再 B 级 M5/M6） \\[2pt]
收尾 & 重做 11/8–11/9 标过的错题（用精编本 / 分享版速查，不翻合订本）；出一张考场时间分配表 \\[2pt]
\end{tabular}
\end{keypoint}

\section{一、早段 · 整卷保温}

\qi{1A}{C 卷（G12）单题回炉 18 题}{材料：61 页 \texttt{每日阅读包\_11-08.pdf} B32，或合订本 p227–232（题面 p227–228 / 答案 p229–232）。\textbf{不掐表}，按单题做，独立写全关键一步。}

\begin{center}\footnotesize
\begin{tabular}{@{}c c p{9.6cm}@{}}
\toprule
题号 & 独立做对 & 卡点 / 备注（错则写原因） \\
C1 & $\square$ & \\
C2 & $\square$ & \\
C3 & $\square$ & \\
C4 & $\square$ & \\
C5 & $\square$ & \\
C6 & $\square$ & \\
C7 & $\square$ & \\
C8 & $\square$ & \\
C9 & $\square$ & \\
C10 & $\square$ & \\
C11 & $\square$ & （重点：该卷第 11 题，压轴形态）\\
C12 & $\square$ & \\
C13 & $\square$ & \\
C14 & $\square$ & \\
C15 & $\square$ & \\
C16 & $\square$ & \\
C17 & $\square$ & \\
C18 & $\square$ & \\
\midrule
\multicolumn{2}{l}{C 卷正确率（做对题数 / 18）} & \blank[3.0cm] \quad 目标 $\ge 80\%$（即 $\ge 15$ 题）\\
\bottomrule
\end{tabular}
\end{center}

\qi{1B}{B 卷 120 min 掐表 ＋ 50 min 订正}{材料：单册 \texttt{模拟卷B.pdf}，或合订本 p221–226。\textbf{严格掐表}，不提前翻答案；做完再统一订正。}

\begin{center}\footnotesize
\begin{tabular}{@{}p{4.2cm} p{2.6cm} p{2.6cm} p{5.2cm}@{}}
\toprule
开始时间 & 结束时间 & 实际用时 & 正确题数 / 18 \\
\blank[3.6cm] & \blank[2.2cm] & \blank[2.2cm] & \blank[3.0cm] \\
\midrule
\multicolumn{4}{l}{18 题结果登记（\yes 对 / \no 错，错则写错因关键词）} \\
B1 \blank[0.8cm] & B4 \blank[0.8cm] & B7 \blank[0.8cm] & B10 \blank[0.8cm] \\
B2 \blank[0.8cm] & B5 \blank[0.8cm] & B8 \blank[0.8cm] & B11 \blank[0.8cm] \\
B3 \blank[0.8cm] & B6 \blank[0.8cm] & B9 \blank[0.8cm] & B12 \blank[0.8cm] \\
B13 \blank[0.8cm] & B14 \blank[0.8cm] & B15 \blank[0.8cm] & B16 \blank[0.8cm] \\
B17 \blank[0.8cm] & B18 \blank[0.8cm] & \multicolumn{2}{l}{订正用时 \blank[2.4cm]} \\
\midrule
\multicolumn{4}{l}{错因归类（对照精编本附录 D 回炉索引，写册/章）} \\
\multicolumn{4}{p{14.6cm}}{\blank[14.0cm]} \\
\bottomrule
\end{tabular}
\end{center}

\section{二、下午 / 晚间 · 定向加练 72 题}

来源：\texttt{讲义\_薄弱点定向加练\_11-09.pdf}（共 84 题，含诊断卷 D1–D12）。\textbf{诊断卷已于 11/9 晚完成（见该册 §0），实测弱项已锁定 M1–M4}，今天直接进入 M 段 72 题，不再重做诊断。按实测顺序 \textbf{A 级四模块优先}；每模块正确率 $<80\%$ 则回对应册重做 3 道同型再往下。做 M1（作图七步）时对照 \texttt{讲义\_函数性态作图\_列表式\_11-09.pdf} 的「性态表五行 = 给分点」写法，列表与图缺一不可。

\begin{center}\footnotesize
\begin{tabular}{@{}c p{5.0cm} c c c c c@{}}
\toprule
模块 & 主题 & 题号 & 题数 & 优先级 & 完成 & 正确率 \\
M1 & 函数性态作图七步（对照列表式讲义） & M1.1–M1.14 & 14 & A-1 & $\square$ & \blank[1.8cm] \\
M2 & 三种量词行 / M 型 / 放缩方向 & M2.1–M2.14 & 14 & A-2 & $\square$ & \blank[1.8cm] \\
M3 & 条件书写 / 六型求导 & M3.1–M3.14 & 14 & A-3 & $\square$ & \blank[1.8cm] \\
M4 & 泰勒：展开阶 / 系数 / 符号 & M4.1–M4.14 & 14 & A-4 & $\square$ & \blank[1.8cm] \\
M5 & 极限计算 / 间断点 & M5.1–M5.8 & 8 & B & $\square$ & \blank[1.8cm] \\
M6 & 新压轴 / 极坐标 & M6.1–M6.8 & 8 & B & $\square$ & \blank[1.8cm] \\
\midrule
\multicolumn{2}{r}{合计} & \multicolumn{2}{c}{72} & \multicolumn{3}{l}{A 级合计 56（占 78\%）} \\
\bottomrule
\end{tabular}
\end{center}

\key{建议顺序：先 M1→M2→M3→M4 把 A 级 56 题刷完（对应你诊断卷权重最高的四弱点），状态好再补 B 级 M5/M6；时间紧则保 M1 与 M4。}

\section{三、验收线 / 放弃线 / 加码}

\begin{warnbox}[不压缩的锚点]
整卷 120 min 掐表、订正 50 min、组末反馈、\textbf{23:00 停笔 / 23:30 睡}——不因题量压缩。
\end{warnbox}

\begin{itemize}\itemsep2pt
\item \textbf{验收线}：B 卷整卷正确率 $\ge 80\%$；定向加练 A 级四模块各 $\ge 80\%$；C 卷 18 题 $\ge 80\%$。
\item \textbf{放弃线}：单题卡 $>8$ min 先跳，组末统一回炉；23:00 到点停笔，未做的顺延到「错题重做」段，不熬夜补。
\item \textbf{加码}：若 18:30 前完成且状态仍好，追加 C 卷第 11 题重点重做，或回炉 11/8–11/9 标过的错题。
\item \textbf{不跳组、不只看答案}：每模块正确率 $<80\%$ 回对应决策卡重做 3 道同型。
\end{itemize}

\section{四、11/11（明天）衔接 · 0 新题}

\begin{keypoint}[考前保温]
\begin{itemize}\itemsep2pt
\item \textbf{0 新题}；只看：精编本附录 A 考前速查卡 + 分享版附录（考前速查 / 17 套必考 / 12 条高频错误）+ 闭卷默写卡（分享版卷首或精编本）。\textbf{合订本 370 页不翻}。
\item 诊断卷 D1–D12 错题重做（昨晚已完成诊断，今天 M 段已覆盖；明早只重做标错的几题，不重新掐表）。
\item 出一张考场时间分配表：选择/填空/计算/证明各分配几分钟，避免考中纠结。
\end{itemize}
\end{keypoint}

\vfill
\noindent{\footnotesize 材料入口（11/9 新产出已并入）：
61 页 \texttt{每日阅读包\_11-08.pdf}（B16–B32）$|$
22 页 \texttt{讲义\_精编本\_11-09.pdf}（方法回炉·轻量主册，8 章方法卡+决策表+章末答案速查）$|$
31 页 \texttt{讲义\_高数I期中\_分享版\_2026-11-09.pdf}（系统通览·考前速查/17 套必考/12 条高频错误附录）$|$
5 页 \texttt{讲义\_函数性态作图\_列表式\_11-09.pdf}（M1 作图七步·性态表=给分点）$|$
370 页 \texttt{高数I期中完整讲义\_合订本\_2026-11-08.pdf}（仅作逐题详解查阅 p162–212，不日常翻）$|$
单册 \texttt{模拟卷B.pdf} $|$ \texttt{讲义\_薄弱点定向加练\_11-09.pdf}（72 题 M 段）。本卡独立 PDF，不改任何既有包的绝对页码。}

\end{document}
